
\cleardoublepage
\chapter{Metodologi Penelitian}
\label{chap:metodologi}

\section{Pendekatan Penelitian}
Pendekatan penelitian yang digunakan dalam studi ini adalah Design Science Research Methodology (DSRM).
DSRM merupakan pendekatan penelitian yang berorientasi pada problem-solving melalui pengembangan artefak yang dapat memberikan solusi terhadap permasalahan nyata dalam suatu organisasi.
Pendekatan ini sangat relevan digunakan dalam penelitian yang bertujuan merancang model, kerangka kerja, atau arsitektur yang bersifat implementatif.

Menurut \textcite{Hevner2004}, Design Science bertujuan menghasilkan artefak yang efektif dan dapat diuji, seperti model, metode, ataupun arsitektur.
Pendekatan ini berupaya menggabungkan proses ilmiah dengan kebutuhan praktis organisasi, sehingga artefak yang dihasilkan tidak hanya memiliki nilai teoritis, tetapi juga memberikan manfaat aplikatif.
Dalam konteks penelitian ini, artefak yang dikembangkan berupa arsitektur transformasi Digital Kemhan RI yang disusun secara terstruktur dengan kerangka TOGAF ADM.

DSRM terdiri dari enam tahapan utama, yaitu: (1) identifikasi masalah dan motivasi, (2) penetapan tujuan solusi, (3) perancangan dan pengembangan artefak, (4) demonstrasi terhadap konteks atau skenario penggunaan, (5) evaluasi artefak, dan (6) komunikasi hasil penelitian. Keenam tahap ini memberikan alur kerja sistematis yang memandu peneliti sejak perumusan masalah hingga penyebaran hasil.

Pendekatan DSRM dipilih karena beberapa alasan utama:
\begin{enumerate}
  \item Penelitian ini berfokus pada pengembangan arsitektur transfomasi digital pada Kementerian Pertahanan yang merupakan artefak desain.
  DSRM menyediakan kerangka yang tepat untuk merancang artefak dengan mempertimbangkan kebutuhan organisasi dan teori yang mendukung.
  \item Transformasi digital Kemhan RI membutuhkan pendekatan yang tidak hanya teoritis, tetapi juga aplikatif.
  DSRM memastikan bahwa artefak yang dihasilkan dapat diterapkan pada konteks nyata organisasi pertahanan.
  \item Melalui tahapan evaluasi, artefak yang dikembangkan dapat dinilai kesesuaiannya oleh para pakar, unit terkait, dan pemangku kepentingan di Kemhan RI.
  Dengan demikian, model yang dihasilkan memiliki validitas yang kuat untuk digunakan dalam desain arsitektur transformasi digital Kementerian Pertahanan.
\end{enumerate}

Dengan demikian, penggunaan DSRM sebagai pendekatan penelitian memberikan landasan metodologis yang kuat dan sesuai untuk menghasilkan
arsitektur transformasi digital Kementerian Pertahanan yang sistematis dan komprehensif.

\section{Tahapan Penelitian Menggunakan DSRM}
Penelitian ini menggunakan Design Science Research Methodology (DSRM) yang terdiri dari enam tahapan utama. 
Setiap tahapan dalam DSRM memberikan struktur penelitian yang sistematis dan berorientasi pada pengembangan artefak yang dapat menyelesaikan permasalahan nyata, dalam hal ini arsitektur transformasi Digital Kemhan RI.
Adapun tahapan penelitian tersebut dijelaskan sebagai berikut:

\subsection{Tahap 1 – Problem Identification and Motivation}
Tahap pertama bertujuan untuk mengidentifikasi permasalahan inti yang menghambat transformasi digital di lingkungan Kemhan RI serta menentukan urgensi penelitian.
Berdasarkan analisis dokumen dan kebijakan, permasalahan utama yang ditemukan adalah sebagai berikut:
\begin{enumerate}
  \item Fragmentasi sistem dan data pertahanan.
  Sistem informasi di Kemhan dan angkatan berjalan sendiri-sendiri, menyebabkan rendahnya interoperabilitas dan integrasi data.
  \item Belum adanya blueprint arsitektur digital Kemhan.
  Kemhan belum memiliki arsitektur SPBE dan arsitektur pertahanan yang terstruktur untuk mendukung modernisasi digital.
  \item Ancaman siber yang meningkat.
  Menurut \textcite{KemhanRI2014} menunjukkan bahwa ancaman siber menjadi tantangan strategis yang belum diimbangi dengan kesiapan sistem dan SDM dalam pertahanan siber.
  \item Keterbatasan SDM digital dan tata kelola TI.
  Kompetensi digital belum merata, dan belum ada mekanisme tata kelola arsitektur yang menyeluruh.
  \item Tuntutan regulasi nasional.
  Berdasarkan \textcite{KementerianKomunikasidanInformatika2018}, \textcite{RepublikIndonesia2023}, \textcite{RI2019}, dan \textcite{RPJPN2025} mewajibkan Kemhan membangun ekosistem digital pertahanan yang terintegrasi.
\end{enumerate}


\subsection{Tahap 2 – Define the Objectives for a Solution}
Ada tahap ini ditetapkan tujuan solusi yang ingin dicapai oleh artefak arsitektur.
Tujuan tersebut disusun berdasarkan permasalahan yang telah diidentifikasi, kebutuhan organisasi, serta mandat regulasi. Adapun tujuan solusi adalah:
\begin{enumerate}
  \item Menghasilkan arsitektur bisnis, data, aplikasi, dan teknologi yang terintegrasi berdasarkan kerangka TOGAF.
  \item Menyediakan blueprint yang dapat digunakan sebagai pedoman transformasi digital Kemhan RI.
  \item Memperkuat tata kelola data pertahanan dan meningkatkan interoperabilitas antar unit organisasi dan matra.
  \item Membangun fondasi penguatan keamanan siber Kemhan.
  \item Mendukung penyusunan roadmap implementasi transformasi digital Kemhan RI 2026–2030.
\end{enumerate}

Tujuan-tujuan ini menjadi dasar dalam pengembangan artefak dan memastikan bahwa solusi yang dihasilkan relevan dengan kondisi aktual dan kebutuhan jangka panjang organisasi.

\subsection{Tahap 3 – Design and Development}
Tabel 3.1 Perancangan Arsitektur Transformasi Digital Kemhan RI

\subsection{Tahap 4 – Demonstration}
Pada tahap ini, artefak yang telah dikembangkan diuji melalui proses demonstrasi untuk menunjukkan bagaimana artefak tersebut menyelesaikan permasalahan yang telah diidentifikasi. Demonstrasi dilakukan melalui:
\begin{enumerate}
  \item Diskusi kelompok terarah (FGD) dengan pakar dan pejabat Kemhan untuk melihat sejauh mana arsitektur yang dirancang dapat diimplementasikan secara praktis.
\end{enumerate}

Demonstrasi bertujuan memastikan bahwa artefak dapat berfungsi secara logis dalam konteks operasional Kemhan RI dan memberikan gambaran implementasi awal sebelum dilakukan evaluasi mendalam.

\subsection{Tahap 5 – Evaluation}
Tahapan ini bertujuan untuk menilai efektivitas artefak arsitektur transformasi Digital Kemhan RI dalam menjawab permasalahan dan mencapai tujuan penelitian. Evaluasi dilakukan dengan dua pendekatan utama:
\begin{enumerate}
  \item Evaluasi Teoritis
  \begin{enumerate}
    \item Menilai kelayakan teknis dan konseptual dari arsitektur transformasi digital yang dirancang,
    termasuk konsistensi dengan prinsip-prinsip arsitektur enterprise (EA) dan keselarasan dengan kerangka kerja TOGAF ADM.
    \item Menilai kesesuaian artefak dengan standar dan kebijakan nasional, seperti arsitektur SPBE Nasional dan regulasi transformasi digital pemerintahan.
    \item Memastikan validitas ilmiah melalui peninjauan oleh pakar akademik dalam bidang sistem informasi, enterprise architecture, dan transformasi digital.
  \end{enumerate}
  \item Evaluasi Praktis
  \begin{enumerate}
    \item Melibatkan pakar dan praktisi dari lingkungan pertahanan, seperti STEI ITB, Ditjen Potensi Pertahanan, Pusdatin Kemhan dan Bainstrahan Kemhan yang relevan untuk memberikan penilaian aplikatif terhadap artefak.
    \item Penilaian dilakukan pada aspek-aspek implementatif, termasuk:
    \begin{enumerate}
      \item Efisiensi Proses
      \item Interoperabilitas sistem pendukung
      \item Keamanan dan kesiapan pengelolaan data
      \item Kelayakan adopsi dalam struktur organisasi Kemhan
    \end{enumerate}
    \item Mengidentifikasi area perbaikan agar arsitektur dapat disesuaikan dengan kondisi operasional Kemhan RI dan kebutuhan transformasi pertahanan.
  \end{enumerate}
\end{enumerate}
Hasil evaluasi memberikan dasar untuk melakukan revisi dan penyempurnaan pada model agar sesuai dengan kebutuhan organisasi dan dinamika lingkungan pertahanan.

\subsection{Tahap 6 – Communication}
Tahap terakhir adalah komunikasi hasil penelitian. Pada tahap ini, artefak dan temuan penelitian disajikan dalam bentuk laporan ilmiah, presentasi kepada pemangku kepentingan, serta publikasi akademik apabila relevan.
Komunikasi bertujuan memastikan bahwa hasil penelitian dapat dipahami, digunakan, dan diadopsi oleh pihak yang berkepentingan.

Dalam konteks penelitian ini, komunikasi melibatkan penyampaian model arsitektur transformasi digital Kemhan RI, dan roadmap implementasi kepada Kemhan RI sebagai pengguna utama.
Dengan demikian, artefak tidak hanya menjadi kontribusi teoretis, tetapi juga memberikan manfaat praktis bagi organisasi.

\section{Jenis dan Sumber Data}
Penelitian ini menggunakan dua jenis data utama, yaitu data primer dan data sekunder, yang dikumpulkan untuk mendukung proses analisis pada setiap tahapan DSRM.
Penggunaan kedua jenis data ini diperlukan agar perancangan arsitektur transformasi digital Kemhan RI memiliki landasan konseptual yang kuat sekaligus relevansi praktis terhadap kondisi aktual di Kementerian Pertahanan RI.
\begin{enumerate}
  \item Data Primer.
  Data primer merupakan data yang diperoleh langsung dari sumber pertama melalui interaksi dengan pemangku kepentingan.
  Data ini diperlukan untuk memahami secara mendalam kondisi aktual sistem informasi, proses bisnis, kebutuhan organisasi, serta tantangan transformasi digital di lingkungan Kemhan. Data primer meliputi:
  \begin{enumerate}
    \item Hasil wawancara dengan pejabat atau staf Kemhan yang memahami sistem informasi, teknologi pertahanan, atau kebijakan SPBE.
    \item Validasi pakar (expert judgement) melalui akademisi atau ahli di bidang enterprise architecture dan transformasi digital pertahanan.
  \end{enumerate}
  Data primer ini sangat penting untuk mengidentifikasi kesenjangan (gap) antara kondisi eksisting dan kebutuhan arsitektur target.
  \item Data Sekunder.
  Data sekunder diperoleh dari dokumen dan sumber referensi yang relevan, antara lain:
  \begin{enumerate}
    \item Peraturan Presiden No 95 Tahun 2018 tentang SPBE.
    \item Peraturan Presiden No 39 Tahun 2019 tentang Satu Data Indonesia.
    \item Peraturan Presiden No 82 Tahun 2023 tentang Percepatan Transformasi Digital.
    \item Rencana Pembangunan Jangka Panjang Nasional (RPJPN) 2025-2045.
    \item Peraturan Menteri Pertahanan No 42 Tahun 2014 tentang Pertahanan Siber.
    \item Undang Undang No 23 Tahun 2019 tentang Pengelolaan Sumber Daya Nasional untuk Pertahanan Negara.
    \item Rencana Strategis Kementerian Pertahanan Tahun 2020-2024.
    \item Jurnal ilmiah, buku, dan publikasi akademik yang membahas transformasi digital, enterprise architecture, kompetensi digital, serta pengembangan SDM di sektor pertahanan.
    \item Referensi internasional, termasuk DoD Digital Modernization Strategy dan penelitian global mengenai digital defense capability, untuk membandingkan praktik terbaik dalam transformasi digital pertahanan.
  \end{enumerate}
  Data sekunder digunakan sebagai landasan teoretis untuk menyusun kerangka konseptual, mendukung analisis gap kompetensi, serta memastikan bahwa arsitektur yang dihasilkan selaras dengan perkembangan global dan kebijakan nasional.
\end{enumerate}

\section{Metode Pengumpulan Data}
Metode pengumpulan data dilakukan melalui beberapa tahapan:
\begin{enumerate}
  \item Studi Literatur.
  Mengkaji teori dasar EA, TOGAF, Smart Defense, dan Transformasi Digital.
  \item Analisis Dokumen Internal.
  Analisis dokumen dilakukan terhadap data dan kebijakan yang dikeluarkan oleh Kemhan RI.
  \item Wawancara Terarah.
  Wawancara dilakukan dengan pejabat atau personel dari unit-unit Kemhan yang relevan, seperti Ditjen Potensi Pertahanan, Pusdatin Kemhan dan Bainstrahan Kemhan. Wawancara ini bertujuan untuk:
  \begin{enumerate}
    \item Menggali informasi mengenai tantangan implementasi transformasi digital.
    \item Mendapatkan perspektif mengenai arah transformasi digital Kemhan RI.
    \item Memvalidasi temuan awal dari analisis dokumen.
  \end{enumerate}
  Metode wawancara semi-terstruktur memungkinkan penggalian informasi yang fleksibel namun tetap terfokus pada isu penelitian.
  \item Focus Group Discussion (FGD).
  FGD digunakan pada tahap evaluasi artefak, sebagai bagian dari pendekatan evaluasi praktis pada DSRM. FGD dilakukan dengan melibatkan:
  \begin{enumerate}
    \item Akademisi
    \item Ditjen Pothan
    \item Pusdatin Kemhan
    \item Bainstrahan Kemhan
  \end{enumerate}
  FGD bertujuan menilai kelayakan arsitektur transformasi digital dari aspek teknis, organisasi, keamanan, dan kesiapan implementasi.
\end{enumerate}













